\documentclass[
  doc,
  floatsintext,
  longtable,
  a4paper,
  nolmodern,
  notxfonts,
  notimes,
  12pt,
  colorlinks=true,linkcolor=blue,citecolor=blue,urlcolor=blue]{apa7}

\usepackage{amsmath}
\usepackage{amssymb}

\geometry{inner=1in, outer=1in}
\fancyhfoffset[LE,RO]{0cm}


\usepackage[bidi=default]{babel}
\babelprovide[main,import]{spanish}


% get rid of language-specific shorthands (see #6817):
\let\LanguageShortHands\languageshorthands
\def\languageshorthands#1{}

\RequirePackage{longtable}
\RequirePackage{threeparttablex}

\makeatletter
\renewcommand{\paragraph}{\@startsection{paragraph}{4}{\parindent}%
	{0\baselineskip \@plus 0.2ex \@minus 0.2ex}%
	{-.5em}%
	{\normalfont\normalsize\bfseries\typesectitle}}

\renewcommand{\subparagraph}[1]{\@startsection{subparagraph}{5}{0.5em}%
	{0\baselineskip \@plus 0.2ex \@minus 0.2ex}%
	{-\z@\relax}%
	{\normalfont\normalsize\bfseries\itshape\hspace{\parindent}{#1}\textit{\addperi}}{\relax}}
\makeatother




\usepackage{longtable, booktabs, multirow, multicol, colortbl, hhline, caption, array, float, xpatch}
\usepackage{subcaption}


\renewcommand\thesubfigure{\Alph{subfigure}}
\setcounter{topnumber}{2}
\setcounter{bottomnumber}{2}
\setcounter{totalnumber}{4}
\renewcommand{\topfraction}{0.85}
\renewcommand{\bottomfraction}{0.85}
\renewcommand{\textfraction}{0.15}
\renewcommand{\floatpagefraction}{0.7}

\usepackage{tcolorbox}
\tcbuselibrary{listings,theorems, breakable, skins}
\usepackage{fontawesome5}

\definecolor{quarto-callout-color}{HTML}{909090}
\definecolor{quarto-callout-note-color}{HTML}{0758E5}
\definecolor{quarto-callout-important-color}{HTML}{CC1914}
\definecolor{quarto-callout-warning-color}{HTML}{EB9113}
\definecolor{quarto-callout-tip-color}{HTML}{00A047}
\definecolor{quarto-callout-caution-color}{HTML}{FC5300}
\definecolor{quarto-callout-color-frame}{HTML}{ACACAC}
\definecolor{quarto-callout-note-color-frame}{HTML}{4582EC}
\definecolor{quarto-callout-important-color-frame}{HTML}{D9534F}
\definecolor{quarto-callout-warning-color-frame}{HTML}{F0AD4E}
\definecolor{quarto-callout-tip-color-frame}{HTML}{02B875}
\definecolor{quarto-callout-caution-color-frame}{HTML}{FD7E14}

%\newlength\Oldarrayrulewidth
%\newlength\Oldtabcolsep


\usepackage{hyperref}



\usepackage{color}
\usepackage{fancyvrb}
\newcommand{\VerbBar}{|}
\newcommand{\VERB}{\Verb[commandchars=\\\{\}]}
\DefineVerbatimEnvironment{Highlighting}{Verbatim}{commandchars=\\\{\}}
% Add ',fontsize=\small' for more characters per line
\usepackage{framed}
\definecolor{shadecolor}{RGB}{241,243,245}
\newenvironment{Shaded}{\begin{snugshade}}{\end{snugshade}}
\newcommand{\AlertTok}[1]{\textcolor[rgb]{0.68,0.00,0.00}{#1}}
\newcommand{\AnnotationTok}[1]{\textcolor[rgb]{0.37,0.37,0.37}{#1}}
\newcommand{\AttributeTok}[1]{\textcolor[rgb]{0.40,0.45,0.13}{#1}}
\newcommand{\BaseNTok}[1]{\textcolor[rgb]{0.68,0.00,0.00}{#1}}
\newcommand{\BuiltInTok}[1]{\textcolor[rgb]{0.00,0.23,0.31}{#1}}
\newcommand{\CharTok}[1]{\textcolor[rgb]{0.13,0.47,0.30}{#1}}
\newcommand{\CommentTok}[1]{\textcolor[rgb]{0.37,0.37,0.37}{#1}}
\newcommand{\CommentVarTok}[1]{\textcolor[rgb]{0.37,0.37,0.37}{\textit{#1}}}
\newcommand{\ConstantTok}[1]{\textcolor[rgb]{0.56,0.35,0.01}{#1}}
\newcommand{\ControlFlowTok}[1]{\textcolor[rgb]{0.00,0.23,0.31}{\textbf{#1}}}
\newcommand{\DataTypeTok}[1]{\textcolor[rgb]{0.68,0.00,0.00}{#1}}
\newcommand{\DecValTok}[1]{\textcolor[rgb]{0.68,0.00,0.00}{#1}}
\newcommand{\DocumentationTok}[1]{\textcolor[rgb]{0.37,0.37,0.37}{\textit{#1}}}
\newcommand{\ErrorTok}[1]{\textcolor[rgb]{0.68,0.00,0.00}{#1}}
\newcommand{\ExtensionTok}[1]{\textcolor[rgb]{0.00,0.23,0.31}{#1}}
\newcommand{\FloatTok}[1]{\textcolor[rgb]{0.68,0.00,0.00}{#1}}
\newcommand{\FunctionTok}[1]{\textcolor[rgb]{0.28,0.35,0.67}{#1}}
\newcommand{\ImportTok}[1]{\textcolor[rgb]{0.00,0.46,0.62}{#1}}
\newcommand{\InformationTok}[1]{\textcolor[rgb]{0.37,0.37,0.37}{#1}}
\newcommand{\KeywordTok}[1]{\textcolor[rgb]{0.00,0.23,0.31}{\textbf{#1}}}
\newcommand{\NormalTok}[1]{\textcolor[rgb]{0.00,0.23,0.31}{#1}}
\newcommand{\OperatorTok}[1]{\textcolor[rgb]{0.37,0.37,0.37}{#1}}
\newcommand{\OtherTok}[1]{\textcolor[rgb]{0.00,0.23,0.31}{#1}}
\newcommand{\PreprocessorTok}[1]{\textcolor[rgb]{0.68,0.00,0.00}{#1}}
\newcommand{\RegionMarkerTok}[1]{\textcolor[rgb]{0.00,0.23,0.31}{#1}}
\newcommand{\SpecialCharTok}[1]{\textcolor[rgb]{0.37,0.37,0.37}{#1}}
\newcommand{\SpecialStringTok}[1]{\textcolor[rgb]{0.13,0.47,0.30}{#1}}
\newcommand{\StringTok}[1]{\textcolor[rgb]{0.13,0.47,0.30}{#1}}
\newcommand{\VariableTok}[1]{\textcolor[rgb]{0.07,0.07,0.07}{#1}}
\newcommand{\VerbatimStringTok}[1]{\textcolor[rgb]{0.13,0.47,0.30}{#1}}
\newcommand{\WarningTok}[1]{\textcolor[rgb]{0.37,0.37,0.37}{\textit{#1}}}

\providecommand{\tightlist}{%
  \setlength{\itemsep}{0pt}\setlength{\parskip}{0pt}}
\usepackage{longtable,booktabs,array}
\newcounter{none} % for unnumbered tables
\usepackage{calc} % for calculating minipage widths
% Correct order of tables after \paragraph or \subparagraph
\usepackage{etoolbox}
\makeatletter
\patchcmd\longtable{\par}{\if@noskipsec\mbox{}\fi\par}{}{}
\makeatother
% Allow footnotes in longtable head/foot
\IfFileExists{footnotehyper.sty}{\usepackage{footnotehyper}}{\usepackage{footnote}}
\makesavenoteenv{longtable}

\usepackage{graphicx}
\makeatletter
\newsavebox\pandoc@box
\newcommand*\pandocbounded[1]{% scales image to fit in text height/width
  \sbox\pandoc@box{#1}%
  \Gscale@div\@tempa{\textheight}{\dimexpr\ht\pandoc@box+\dp\pandoc@box\relax}%
  \Gscale@div\@tempb{\linewidth}{\wd\pandoc@box}%
  \ifdim\@tempb\p@<\@tempa\p@\let\@tempa\@tempb\fi% select the smaller of both
  \ifdim\@tempa\p@<\p@\scalebox{\@tempa}{\usebox\pandoc@box}%
  \else\usebox{\pandoc@box}%
  \fi%
}
% Set default figure placement to htbp
\def\fps@figure{htbp}
\makeatother







\usepackage{newtx}

\defaultfontfeatures{Scale=MatchLowercase}
\defaultfontfeatures[\rmfamily]{Ligatures=TeX,Scale=1}





\title{Rebanado de cadenas}


\shorttitle{REBANADO DE CADENAS}


\usepackage{etoolbox}



\ccoppy{\textcopyright~2026}



\author{Edison Achalma}



\affiliation{
{Escuela Profesional de Economía, Universidad Nacional de San Cristóbal
de Huamanga}}




\leftheader{Achalma}

\date{2026-10-20}


\abstract{Este abstract será actualizado una vez que se complete el
contenido final del artículo. }



\authornote{\par{\addORCIDlink{Edison Achalma}{0000-0001-6996-3364}} 
\par{ }
\par{   El autor no tiene conflictos de interés que revelar.    Los
roles de autor se clasificaron utilizando la taxonomía de roles de
colaborador (CRediT; https://credit.niso.org/) de la siguiente
manera:  Edison Achalma:   conceptualización, metodología, análisis
formal, investigación, recursos, curación de
datos, redacción, visualización, supervisión, administración del
proyecto}
\par{La correspondencia relativa a este artículo debe dirigirse a Edison
Achalma, Escuela Profesional de Economía, Universidad Nacional de San
Cristóbal de Huamanga, Portal Independencia N°
57, Ayacucho, AYA 5001, Perú, Email: \href{mailto:elmer.achalma.09@unsch.edu.pe}{elmer.achalma.09@unsch.edu.pe}}
}

\makeatletter
\let\endoldlt\endlongtable
\def\endlongtable{
\hline
\endoldlt
}
\makeatother

\urlstyle{same}



\hypersetup{
  pdftitle={Rebanado de cadenas},
  pdfauthor={Autor}
}
\makeatletter
\@ifpackageloaded{caption}{}{\usepackage{caption}}
\AtBeginDocument{%
\ifdefined\contentsname
  \renewcommand*\contentsname{Tabla de contenidos}
\else
  \newcommand\contentsname{Tabla de contenidos}
\fi
\ifdefined\listfigurename
  \renewcommand*\listfigurename{Lista de Figuras}
\else
  \newcommand\listfigurename{Lista de Figuras}
\fi
\ifdefined\listtablename
  \renewcommand*\listtablename{Lista de Tablas}
\else
  \newcommand\listtablename{Lista de Tablas}
\fi
\ifdefined\figurename
  \renewcommand*\figurename{Figura}
\else
  \newcommand\figurename{Figura}
\fi
\ifdefined\tablename
  \renewcommand*\tablename{Tabla}
\else
  \newcommand\tablename{Tabla}
\fi
}
\@ifpackageloaded{float}{}{\usepackage{float}}
\floatstyle{ruled}
\@ifundefined{c@chapter}{\newfloat{codelisting}{h}{lop}}{\newfloat{codelisting}{h}{lop}[chapter]}
\floatname{codelisting}{Listado}
\newcommand*\listoflistings{\listof{codelisting}{Lista de Listados}}
\makeatother
\makeatletter
\makeatother
\makeatletter
\@ifpackageloaded{caption}{}{\usepackage{caption}}
\@ifpackageloaded{subcaption}{}{\usepackage{subcaption}}
\makeatother
\makeatletter
\@ifpackageloaded{fontawesome5}{}{\usepackage{fontawesome5}}
\makeatother

% From https://tex.stackexchange.com/a/645996/211326
%%% apa7 doesn't want to add appendix section titles in the toc
%%% let's make it do it
\makeatletter
\xpatchcmd{\appendix}
  {\par}
  {\addcontentsline{toc}{section}{\@currentlabelname}\par}
  {}{}
\makeatother

%% Disable longtable counter
%% https://tex.stackexchange.com/a/248395/211326

\usepackage{etoolbox}

\makeatletter
\patchcmd{\LT@caption}
  {\bgroup}
  {\bgroup\global\LTpatch@captiontrue}
  {}{}
\patchcmd{\longtable}
  {\par}
  {\par\global\LTpatch@captionfalse}
  {}{}
\apptocmd{\endlongtable}
  {\ifLTpatch@caption\else\addtocounter{table}{-1}\fi}
  {}{}
\newif\ifLTpatch@caption
\makeatother

\begin{document}

\maketitle


\hypertarget{toc}{}
\tableofcontents
\newpage
\section[Introduction]{Rebanado de cadenas}

\setcounter{secnumdepth}{3}

\setlength\LTleft{0pt}




\subsection{🎯 Motivación}\label{motivaciuxf3n}

A menudo necesitas trabajar con partes de un texto: extraer el nombre de
usuario de un email, obtener la extensión de un archivo, o simplemente
tomar una porción de un mensaje. El slicing (rebanado) te permite hacer
esto de manera elegante y eficiente, sin necesidad de bucles o lógica
compleja.

\subsection{💡 Intuición}\label{intuiciuxf3n}

Imagina que tienes una fila de personas (string) y quieres seleccionar a
las personas de las posiciones 2 a 5. El slicing es como decir: ``dame
desde la posición 2 hasta la posición 5''. En Python, \texttt{{[}2:5{]}}
significa: empieza en el índice 2 y \textbf{detente antes} del índice 5.
Es como un cuchillo que corta exactamente donde le indicas.

\subsection{🧠 Modelo mental}\label{modelo-mental}

La sintaxis de slicing es \texttt{{[}inicio:fin:paso{]}}:

\begin{itemize}
\tightlist
\item
  \textbf{\texttt{inicio}} (opcional): dónde empezar (incluido). Por
  defecto es el principio.
\item
  \textbf{\texttt{fin}} (opcional): dónde parar (no incluido). Por
  defecto es el final.
\item
  \textbf{\texttt{paso}} (opcional): cada cuántos caracteres tomar. Por
  defecto es 1.
\end{itemize}

La regla clave: \textbf{el índice final es exclusivo}, es decir, no se
incluye en el resultado.

\subsection{📐 Desarrollo}\label{desarrollo}

\subsubsection{1. Sintaxis básica}\label{sintaxis-buxe1sica}

\begin{Shaded}
\begin{Highlighting}[]
\NormalTok{texto }\OperatorTok{=} \StringTok{"Hello, World!"}

\CommentTok{\# [inicio:fin] {-} fin no incluido}
\BuiltInTok{print}\NormalTok{(texto[}\DecValTok{2}\NormalTok{:}\DecValTok{5}\NormalTok{])   }\CommentTok{\# "llo" (índices 2, 3, 4)}
\BuiltInTok{print}\NormalTok{(texto[}\DecValTok{7}\NormalTok{:}\DecValTok{12}\NormalTok{])  }\CommentTok{\# "World" (índices 7{-}11)}
\end{Highlighting}
\end{Shaded}

\subsubsection{2. Slice desde el inicio}\label{slice-desde-el-inicio}

Omite el inicio para empezar desde el primer carácter:

\begin{Shaded}
\begin{Highlighting}[]
\NormalTok{texto }\OperatorTok{=} \StringTok{"Hello, World!"}
\BuiltInTok{print}\NormalTok{(texto[:}\DecValTok{5}\NormalTok{])   }\CommentTok{\# "Hello" (índices 0{-}4)}
\BuiltInTok{print}\NormalTok{(texto[:}\DecValTok{7}\NormalTok{])   }\CommentTok{\# "Hello, " (índices 0{-}6)}
\end{Highlighting}
\end{Shaded}

\subsubsection{3. Slice hasta el final}\label{slice-hasta-el-final}

Omite el fin para ir hasta el último carácter:

\begin{Shaded}
\begin{Highlighting}[]
\NormalTok{texto }\OperatorTok{=} \StringTok{"Hello, World!"}
\BuiltInTok{print}\NormalTok{(texto[}\DecValTok{7}\NormalTok{:])   }\CommentTok{\# "World!" (índices 7{-}12)}
\BuiltInTok{print}\NormalTok{(texto[}\DecValTok{2}\NormalTok{:])   }\CommentTok{\# "llo, World!" (índices 2{-}12)}
\end{Highlighting}
\end{Shaded}

\subsubsection{4. Slice completo}\label{slice-completo}

\begin{Shaded}
\begin{Highlighting}[]
\NormalTok{texto }\OperatorTok{=} \StringTok{"Hello, World!"}
\BuiltInTok{print}\NormalTok{(texto[:])    }\CommentTok{\# "Hello, World!" (copia completa)}
\end{Highlighting}
\end{Shaded}

\subsubsection{5. Índices negativos}\label{uxedndices-negativos}

Empieza desde el final (último carácter es -1):

\begin{Shaded}
\begin{Highlighting}[]
\NormalTok{texto }\OperatorTok{=} \StringTok{"Hello, World!"}

\CommentTok{\# Últimos 5 caracteres}
\BuiltInTok{print}\NormalTok{(texto[}\OperatorTok{{-}}\DecValTok{5}\NormalTok{:])   }\CommentTok{\# "orld!"}
\BuiltInTok{print}\NormalTok{(texto[}\OperatorTok{{-}}\DecValTok{6}\NormalTok{:}\OperatorTok{{-}}\DecValTok{1}\NormalTok{]) }\CommentTok{\# "World" (desde {-}6 hasta antes de {-}1)}
\BuiltInTok{print}\NormalTok{(texto[}\OperatorTok{{-}}\DecValTok{5}\NormalTok{:}\OperatorTok{{-}}\DecValTok{2}\NormalTok{]) }\CommentTok{\# "orl" (desde {-}5 hasta antes de {-}2)}
\end{Highlighting}
\end{Shaded}

\subsubsection{\texorpdfstring{6. Slice con paso
(\texttt{{[}inicio:fin:paso{]}})}{6. Slice con paso ({[}inicio:fin:paso{]})}}\label{slice-con-paso-iniciofinpaso}

Controlas el salto entre caracteres:

\begin{Shaded}
\begin{Highlighting}[]
\NormalTok{texto }\OperatorTok{=} \StringTok{"Python"}

\CommentTok{\# Cada 2 caracteres}
\BuiltInTok{print}\NormalTok{(texto[::}\DecValTok{2}\NormalTok{])   }\CommentTok{\# "Pto" (índices 0, 2, 4)}
\BuiltInTok{print}\NormalTok{(texto[}\DecValTok{1}\NormalTok{::}\DecValTok{2}\NormalTok{])  }\CommentTok{\# "yhn" (índices 1, 3, 5)}

\CommentTok{\# Invertir string (paso {-}1)}
\BuiltInTok{print}\NormalTok{(texto[::}\OperatorTok{{-}}\DecValTok{1}\NormalTok{])  }\CommentTok{\# "nohtyP"}

\CommentTok{\# Desde el final con paso {-}2}
\BuiltInTok{print}\NormalTok{(texto[::}\OperatorTok{{-}}\DecValTok{2}\NormalTok{])  }\CommentTok{\# "nhy"}
\end{Highlighting}
\end{Shaded}

\subsubsection{7. Tabla de ejemplos de
slicing}\label{tabla-de-ejemplos-de-slicing}

{\def\LTcaptype{none} % do not increment counter
\begin{longtable}[]{@{}
  >{\raggedright\arraybackslash}p{(\linewidth - 4\tabcolsep) * \real{0.1231}}
  >{\raggedright\arraybackslash}p{(\linewidth - 4\tabcolsep) * \real{0.3846}}
  >{\raggedright\arraybackslash}p{(\linewidth - 4\tabcolsep) * \real{0.4923}}@{}}
\toprule\noalign{}
\begin{minipage}[b]{\linewidth}\raggedright
Slicing
\end{minipage} & \begin{minipage}[b]{\linewidth}\raggedright
Resultado (para ``Python'')
\end{minipage} & \begin{minipage}[b]{\linewidth}\raggedright
Explicación
\end{minipage} \\
\midrule\noalign{}
\endhead
\bottomrule\noalign{}
\endlastfoot
\texttt{{[}0:3{]}} & \texttt{"Pyt"} & Posiciones 0, 1, 2 \\
\texttt{{[}2:5{]}} & \texttt{"tho"} & Posiciones 2, 3, 4 \\
\texttt{{[}:4{]}} & \texttt{"Pyth"} & Desde el inicio hasta posición
3 \\
\texttt{{[}3:{]}} & \texttt{"hon"} & Desde posición 3 hasta el final \\
\texttt{{[}-3:{]}} & \texttt{"hon"} & Últimos 3 caracteres \\
\texttt{{[}:-2{]}} & \texttt{"Pyth"} & Hasta antes de los últimos 2 \\
\texttt{{[}::2{]}} & \texttt{"Pto"} & Caracteres en posiciones pares \\
\texttt{{[}::-1{]}} & \texttt{"nohtyP"} & String invertido \\
\end{longtable}
}

\subsection{🔬 Ejemplos}\label{ejemplos}

\textbf{Mínimo ---} Slicing básico:

\begin{Shaded}
\begin{Highlighting}[]
\NormalTok{fruta }\OperatorTok{=} \StringTok{"Manzana"}
\BuiltInTok{print}\NormalTok{(fruta[}\DecValTok{0}\NormalTok{:}\DecValTok{3}\NormalTok{])  }\CommentTok{\# "Man"}
\BuiltInTok{print}\NormalTok{(fruta[}\DecValTok{3}\NormalTok{:}\DecValTok{6}\NormalTok{])  }\CommentTok{\# "ana"}
\BuiltInTok{print}\NormalTok{(fruta[}\DecValTok{2}\NormalTok{:}\DecValTok{5}\NormalTok{])  }\CommentTok{\# "nza"}
\end{Highlighting}
\end{Shaded}

\textbf{Progresivo ---} Extracción de datos:

\begin{Shaded}
\begin{Highlighting}[]
\NormalTok{email }\OperatorTok{=} \StringTok{"usuario@gmail.com"}

\CommentTok{\# Extraer nombre de usuario (hasta el @)}
\NormalTok{nombre\_usuario }\OperatorTok{=}\NormalTok{ email[:email.index(}\StringTok{"@"}\NormalTok{)]}
\BuiltInTok{print}\NormalTok{(nombre\_usuario)  }\CommentTok{\# "usuario"}

\CommentTok{\# Extraer dominio (después del @)}
\NormalTok{dominio }\OperatorTok{=}\NormalTok{ email[email.index(}\StringTok{"@"}\NormalTok{) }\OperatorTok{+} \DecValTok{1}\NormalTok{:]}
\BuiltInTok{print}\NormalTok{(dominio)  }\CommentTok{\# "gmail.com"}

\CommentTok{\# Extraer extensión (últimos 3 caracteres)}
\NormalTok{extension }\OperatorTok{=}\NormalTok{ email[}\OperatorTok{{-}}\DecValTok{3}\NormalTok{:]}
\BuiltInTok{print}\NormalTok{(extension)  }\CommentTok{\# "com"}
\end{Highlighting}
\end{Shaded}

\textbf{Real ---} Procesamiento de nombres de archivo:

\begin{Shaded}
\begin{Highlighting}[]
\NormalTok{archivo }\OperatorTok{=} \StringTok{"documento\_2024.pdf"}

\CommentTok{\# Obtener extensión}
\NormalTok{extension }\OperatorTok{=}\NormalTok{ archivo.split(}\StringTok{"."}\NormalTok{)[}\OperatorTok{{-}}\DecValTok{1}\NormalTok{]  }\CommentTok{\# "pdf"}
\CommentTok{\# Con slicing:}
\NormalTok{punto }\OperatorTok{=}\NormalTok{ archivo.index(}\StringTok{"."}\NormalTok{)}
\NormalTok{extension\_slice }\OperatorTok{=}\NormalTok{ archivo[punto }\OperatorTok{+} \DecValTok{1}\NormalTok{:]  }\CommentTok{\# "pdf"}

\CommentTok{\# Obtener nombre sin extensión}
\NormalTok{nombre\_base }\OperatorTok{=}\NormalTok{ archivo[:archivo.index(}\StringTok{"."}\NormalTok{)]  }\CommentTok{\# "documento\_2024"}

\CommentTok{\# Año (últimos 4 dígitos antes del punto)}
\NormalTok{año }\OperatorTok{=}\NormalTok{ archivo[}\OperatorTok{{-}}\DecValTok{8}\NormalTok{:}\OperatorTok{{-}}\DecValTok{4}\NormalTok{]  }\CommentTok{\# "2024"}

\BuiltInTok{print}\NormalTok{(}\SpecialStringTok{f"Archivo: }\SpecialCharTok{\{}\NormalTok{nombre\_base}\SpecialCharTok{\}}\SpecialStringTok{"}\NormalTok{)}
\BuiltInTok{print}\NormalTok{(}\SpecialStringTok{f"Extensión: }\SpecialCharTok{\{}\NormalTok{extension\_slice}\SpecialCharTok{\}}\SpecialStringTok{"}\NormalTok{)}
\BuiltInTok{print}\NormalTok{(}\SpecialStringTok{f"Año: }\SpecialCharTok{\{}\NormalTok{a}\SpecialCharTok{ñ}\NormalTok{o}\SpecialCharTok{\}}\SpecialStringTok{"}\NormalTok{)}
\end{Highlighting}
\end{Shaded}

\subsection{⚠️ Errores comunes}\label{errores-comunes}

{\def\LTcaptype{none} % do not increment counter
\begin{longtable}[]{@{}
  >{\raggedright\arraybackslash}p{(\linewidth - 4\tabcolsep) * \real{0.2752}}
  >{\raggedright\arraybackslash}p{(\linewidth - 4\tabcolsep) * \real{0.3578}}
  >{\raggedright\arraybackslash}p{(\linewidth - 4\tabcolsep) * \real{0.3670}}@{}}
\toprule\noalign{}
\begin{minipage}[b]{\linewidth}\raggedright
Error
\end{minipage} & \begin{minipage}[b]{\linewidth}\raggedright
Síntoma
\end{minipage} & \begin{minipage}[b]{\linewidth}\raggedright
Solución
\end{minipage} \\
\midrule\noalign{}
\endhead
\bottomrule\noalign{}
\endlastfoot
\textbf{Índice fuera de rango} &
\texttt{IndexError:\ string\ index\ out\ of\ range} & Verifica que el
índice sea \textless{} len(string) \\
\textbf{Confundir \texttt{{[}inicio:fin{]}}} & No incluir el último
carácter & Recuerda: \texttt{fin} no se incluye \\
\textbf{Olvidar \texttt{:}} & Intenta acceder a un solo carácter & Usa
\texttt{{[}inicio:fin{]}} no \texttt{{[}inicio,\ fin{]}} \\
\textbf{Índices negativos confusos} & Resultados inesperados & Recuerda:
-1 es el último carácter \\
\end{longtable}
}

\subsection{✅ Buenas prácticas}\label{buenas-pruxe1cticas}

\begin{enumerate}
\def\labelenumi{\arabic{enumi}.}
\item
  \textbf{Usa índices negativos para ``desde el final''}:

\begin{Shaded}
\begin{Highlighting}[]
\CommentTok{\# ✅ Obtener últimos 4 caracteres}
\NormalTok{ultimos\_4 }\OperatorTok{=}\NormalTok{ texto[}\OperatorTok{{-}}\DecValTok{4}\NormalTok{:]}

\CommentTok{\# ❌ Menos claro}
\NormalTok{ultimos\_4 }\OperatorTok{=}\NormalTok{ texto[}\BuiltInTok{len}\NormalTok{(texto)}\OperatorTok{{-}}\DecValTok{4}\NormalTok{:]}
\end{Highlighting}
\end{Shaded}
\item
  \textbf{Combina slicing con \texttt{len()}} para cortar dinámicamente:

\begin{Shaded}
\begin{Highlighting}[]
\NormalTok{texto }\OperatorTok{=} \StringTok{"Hola mundo"}
\CommentTok{\# Primeros 3}
\NormalTok{texto[:}\DecValTok{3}\NormalTok{]}
\CommentTok{\# Últimos 3}
\NormalTok{texto[}\OperatorTok{{-}}\DecValTok{3}\NormalTok{:]}
\CommentTok{\# Mitad}
\NormalTok{mitad }\OperatorTok{=} \BuiltInTok{len}\NormalTok{(texto) }\OperatorTok{//} \DecValTok{2}
\NormalTok{texto[:mitad]  }\CommentTok{\# Primera mitad}
\NormalTok{texto[mitad:]  }\CommentTok{\# Segunda mitad}
\end{Highlighting}
\end{Shaded}
\item
  \textbf{Usa \texttt{{[}::-1{]}} para invertir strings}:

\begin{Shaded}
\begin{Highlighting}[]
\NormalTok{palabra }\OperatorTok{=} \StringTok{"Python"}
\BuiltInTok{print}\NormalTok{(palabra[::}\OperatorTok{{-}}\DecValTok{1}\NormalTok{])  }\CommentTok{\# "nohtyP"}
\end{Highlighting}
\end{Shaded}
\item
  \textbf{Verifica que los índices sean válidos} antes de hacer slicing
  si no estás seguro:

\begin{Shaded}
\begin{Highlighting}[]
\NormalTok{texto }\OperatorTok{=} \StringTok{"abc"}
\CommentTok{\# Slice con índices fuera de rango es seguro (no da error)}
\BuiltInTok{print}\NormalTok{(texto[}\DecValTok{0}\NormalTok{:}\DecValTok{10}\NormalTok{])  }\CommentTok{\# "abc" (no da error)}
\CommentTok{\# Pero acceso directo sí da error}
\CommentTok{\# print(texto[5])  \# ❌ IndexError}
\end{Highlighting}
\end{Shaded}
\item
  \textbf{Usa slicing para copiar strings}:

\begin{Shaded}
\begin{Highlighting}[]
\NormalTok{original }\OperatorTok{=} \StringTok{"Hola"}
\NormalTok{copia }\OperatorTok{=}\NormalTok{ original[:]  }\CommentTok{\# Crea una copia}
\end{Highlighting}
\end{Shaded}
\end{enumerate}

\subsection{📝 Ejercicio}\label{ejercicio}

\textbf{Intermedio} --- Manipulación de nombres de archivo

Escribe un programa que:

\begin{enumerate}
\def\labelenumi{\arabic{enumi}.}
\tightlist
\item
  Tenga un string \texttt{ruta\_archivo} =
  ``C:/Users/Ana/Documents/proyecto\_final\_v2.py''
\item
  Extraiga y muestre:

  \begin{itemize}
  \tightlist
  \item
    El nombre del archivo (última parte después de /)
  \item
    La extensión del archivo
  \item
    El nombre base (sin extensión)
  \item
    La carpeta donde está el archivo (todo antes del último /)
  \item
    Una versión con solo las iniciales de cada carpeta:
    C:/U/A/D/proyecto\_final\_v2.py
  \end{itemize}
\end{enumerate}

\textbf{Pista:} usa \texttt{rfind()} para encontrar la última ocurrencia
de un carácter, y luego slicing.

\begin{center}\rule{0.5\linewidth}{0.5pt}\end{center}

\textbf{Solución (no mires hasta haberlo intentado):}

\begin{Shaded}
\begin{Highlighting}[]
\NormalTok{ruta\_archivo }\OperatorTok{=} \StringTok{"C:/Users/Ana/Documents/proyecto\_final\_v2.py"}

\CommentTok{\# 1. Nombre del archivo (último segmento)}
\NormalTok{ultimo\_slash }\OperatorTok{=}\NormalTok{ ruta\_archivo.rfind(}\StringTok{"/"}\NormalTok{)}
\NormalTok{nombre\_archivo }\OperatorTok{=}\NormalTok{ ruta\_archivo[ultimo\_slash }\OperatorTok{+} \DecValTok{1}\NormalTok{:]}
\BuiltInTok{print}\NormalTok{(}\SpecialStringTok{f"Archivo: }\SpecialCharTok{\{}\NormalTok{nombre\_archivo}\SpecialCharTok{\}}\SpecialStringTok{"}\NormalTok{)  }\CommentTok{\# "proyecto\_final\_v2.py"}

\CommentTok{\# 2. Extensión}
\NormalTok{punto }\OperatorTok{=}\NormalTok{ nombre\_archivo.rfind(}\StringTok{"."}\NormalTok{)}
\NormalTok{extension }\OperatorTok{=}\NormalTok{ nombre\_archivo[punto }\OperatorTok{+} \DecValTok{1}\NormalTok{:]}
\BuiltInTok{print}\NormalTok{(}\SpecialStringTok{f"Extensión: }\SpecialCharTok{\{}\NormalTok{extension}\SpecialCharTok{\}}\SpecialStringTok{"}\NormalTok{)  }\CommentTok{\# "py"}

\CommentTok{\# 3. Nombre base (sin extensión)}
\NormalTok{nombre\_base }\OperatorTok{=}\NormalTok{ nombre\_archivo[:punto]}
\BuiltInTok{print}\NormalTok{(}\SpecialStringTok{f"Nombre base: }\SpecialCharTok{\{}\NormalTok{nombre\_base}\SpecialCharTok{\}}\SpecialStringTok{"}\NormalTok{)  }\CommentTok{\# "proyecto\_final\_v2"}

\CommentTok{\# 4. Carpeta (todo antes del último /)}
\NormalTok{carpeta }\OperatorTok{=}\NormalTok{ ruta\_archivo[:ultimo\_slash]}
\BuiltInTok{print}\NormalTok{(}\SpecialStringTok{f"Carpeta: }\SpecialCharTok{\{}\NormalTok{carpeta}\SpecialCharTok{\}}\SpecialStringTok{"}\NormalTok{)  }\CommentTok{\# "C:/Users/Ana/Documents"}

\CommentTok{\# 5. Iniciales de carpetas}
\NormalTok{partes }\OperatorTok{=}\NormalTok{ ruta\_archivo.split(}\StringTok{"/"}\NormalTok{)}
\NormalTok{iniciales }\OperatorTok{=} \StringTok{"/"}\NormalTok{.join([p[:}\DecValTok{1}\NormalTok{] }\ControlFlowTok{if}\NormalTok{ p }\ControlFlowTok{else} \StringTok{""} \ControlFlowTok{for}\NormalTok{ p }\KeywordTok{in}\NormalTok{ partes[:}\OperatorTok{{-}}\DecValTok{1}\NormalTok{]]) }\OperatorTok{+} \StringTok{"/"} \OperatorTok{+}\NormalTok{ partes[}\OperatorTok{{-}}\DecValTok{1}\NormalTok{]}
\BuiltInTok{print}\NormalTok{(}\SpecialStringTok{f"Con iniciales: }\SpecialCharTok{\{}\NormalTok{iniciales}\SpecialCharTok{\}}\SpecialStringTok{"}\NormalTok{)  }\CommentTok{\# "C:/U/A/D/proyecto\_final\_v2.py"}
\end{Highlighting}
\end{Shaded}

\subsection{🗂️ Síntesis}\label{suxedntesis}

\begin{itemize}
\tightlist
\item
  \textbf{Sintaxis}: \texttt{{[}inicio:fin:paso{]}} donde \texttt{fin}
  no se incluye.
\item
  \textbf{Inicio por defecto}: \texttt{{[}:fin{]}} empieza desde el
  principio.
\item
  \textbf{Fin por defecto}: \texttt{{[}inicio:{]}} va hasta el final.
\item
  \textbf{Índices negativos}: cuentan desde el final (-1 es el último).
\item
  \textbf{Paso}: controla el salto entre caracteres. \texttt{{[}::-1{]}}
  invierte.
\item
  \textbf{Slicing seguro}: índices fuera de rango no dan error (a
  diferencia de acceso directo).
\item
  \textbf{Strings inmutables}: el slicing devuelve un nuevo string, no
  modifica el original.
\end{itemize}

\subsection{🔗 Qué sigue}\label{quuxe9-sigue}

\begin{itemize}
\tightlist
\item
  \textbf{Relacionado:} métodos de string, indexación, f-strings
\item
  \textbf{Siguiente tema:} Métodos de strings --- aprende todas las
  herramientas para manipular texto.
\end{itemize}

\section{Publicaciones Similares}\label{publicaciones-similares}

Si te interesó este artículo, te recomendamos que explores otros blogs y
recursos relacionados que pueden ampliar tus conocimientos. Aquí te dejo
algunas sugerencias:

\begin{enumerate}
\def\labelenumi{\arabic{enumi}.}
\tightlist
\item
  \href{https://numerus-scriptum.netlify.app/python/2020-06-19-instalacion-de-anaconda/index.pdf}{\faIcon{file-pdf}}
  \href{https://numerus-scriptum.netlify.app/python/2020-06-19-instalacion-de-anaconda}{Instalacion
  De Anaconda}
\item
  \href{https://numerus-scriptum.netlify.app/python/2020-06-20-configurar-entorno-virtual-python-anaconda/index.pdf}{\faIcon{file-pdf}}
  \href{https://numerus-scriptum.netlify.app/python/2020-06-20-configurar-entorno-virtual-python-anaconda}{Configurar
  Entorno Virtual Python Anaconda}
\item
  \href{https://numerus-scriptum.netlify.app/python/2021-04-17-01-introducion-a-la-programacion-con-python/index.pdf}{\faIcon{file-pdf}}
  \href{https://numerus-scriptum.netlify.app/python/2021-04-17-01-introducion-a-la-programacion-con-python}{01
  Introducion A La Programacion Con Python}
\item
  \href{https://numerus-scriptum.netlify.app/python/2021-05-31-02-variables-expresiones-y-statements-con-python/index.pdf}{\faIcon{file-pdf}}
  \href{https://numerus-scriptum.netlify.app/python/2021-05-31-02-variables-expresiones-y-statements-con-python}{02
  Variables Expresiones Y Statements Con Python}
\item
  \href{https://numerus-scriptum.netlify.app/python/2021-06-07-03-objetos-de-python/index.pdf}{\faIcon{file-pdf}}
  \href{https://numerus-scriptum.netlify.app/python/2021-06-07-03-objetos-de-python}{03
  Objetos De Python}
\item
  \href{https://numerus-scriptum.netlify.app/python/2021-06-14-04-ejecucion-condicional-con-python/index.pdf}{\faIcon{file-pdf}}
  \href{https://numerus-scriptum.netlify.app/python/2021-06-14-04-ejecucion-condicional-con-python}{04
  Ejecucion Condicional Con Python}
\item
  \href{https://numerus-scriptum.netlify.app/python/2021-06-21-05-iteraciones-con-python/index.pdf}{\faIcon{file-pdf}}
  \href{https://numerus-scriptum.netlify.app/python/2021-06-21-05-iteraciones-con-python}{05
  Iteraciones Con Python}
\item
  \href{https://numerus-scriptum.netlify.app/python/2021-08-16-06-funciones-con-python/index.pdf}{\faIcon{file-pdf}}
  \href{https://numerus-scriptum.netlify.app/python/2021-08-16-06-funciones-con-python}{06
  Funciones Con Python}
\item
  \href{https://numerus-scriptum.netlify.app/python/2021-08-23-07-dataframes-con-python/index.pdf}{\faIcon{file-pdf}}
  \href{https://numerus-scriptum.netlify.app/python/2021-08-23-07-dataframes-con-python}{07
  Dataframes Con Python}
\item
  \href{https://numerus-scriptum.netlify.app/python/2021-11-29-08-prediccion-y-metrica-de-performance-con-python/index.pdf}{\faIcon{file-pdf}}
  \href{https://numerus-scriptum.netlify.app/python/2021-11-29-08-prediccion-y-metrica-de-performance-con-python}{08
  Prediccion Y Metrica De Performance Con Python}
\item
  \href{https://numerus-scriptum.netlify.app/python/2021-12-06-09-metodos-de-machine-learning-para-clasificacion-con-python/index.pdf}{\faIcon{file-pdf}}
  \href{https://numerus-scriptum.netlify.app/python/2021-12-06-09-metodos-de-machine-learning-para-clasificacion-con-python}{09
  Metodos De Machine Learning Para Clasificacion Con Python}
\item
  \href{https://numerus-scriptum.netlify.app/python/2021-12-13-10-metodos-de-machine-learning-para-regresion-con-python/index.pdf}{\faIcon{file-pdf}}
  \href{https://numerus-scriptum.netlify.app/python/2021-12-13-10-metodos-de-machine-learning-para-regresion-con-python}{10
  Metodos De Machine Learning Para Regresion Con Python}
\item
  \href{https://numerus-scriptum.netlify.app/python/2025-05-10-visualizacion-de-datos-con-python/index.pdf}{\faIcon{file-pdf}}
  \href{https://numerus-scriptum.netlify.app/python/2025-05-10-visualizacion-de-datos-con-python}{Visualizacion
  De Datos Con Python}
\end{enumerate}

Esperamos que encuentres estas publicaciones igualmente interesantes y
útiles. ¡Disfruta de la lectura!






\end{document}
